%% P05 \dash{} Iloczyn skalarny, normy i rzutowanie
%%
%% Zalążek. Poniższy opis jest kontraktem tego programu: co ma uzasadnić i co
%% ma dać czytelnikowi. Napisanie programu oznacza usunięcie bloku
%% \programstub{} i zastąpienie go programem. Jeśli po przerobieniu materiału
%% opis okaże się błędny, popraw go i odnotuj sprzeczność w CLAUDE.md w sekcji
%% *Resolved questions*.

\program{Iloczyn skalarny, normy i rzutowanie}
\label{prog:P05}

\programstub{%
\textit{[Opis do przetłumaczenia \dash{} patrz wydanie angielskie.]}

Argument: an inner product is the only thing that gives a vector space
angles and lengths; a projection is the closest point in a subspace;
different norms measure different things and disagree about which vector is
bigger. Payoff: dot product, cosine and Euclidean distance compared on the
same data, with the case where they rank differently worked out \dash{} the
thing a vector database forces you to choose and nobody explains; why
normalising embeddings makes cosine and dot product the same query and what
you lose; L1 against L2 as a shape rather than a slogan, and why the L1
ball's corners are the whole of the sparsity argument. Also the fact the
book cashes in twice later: in high dimension, two random vectors are almost
always nearly orthogonal, measured rather than asserted.

\medskip
\textbf{Planowana długość:} 60 ramek.
}
